\documentclass[10pt,a4paper]{amsart}
\usepackage[margin=19mm]{geometry}
\usepackage{amsmath,amssymb,mathtools}
\usepackage[scr=rsfs,cal=euler]{mathalfa}
\usepackage{xfrac}
\usepackage[unicode,hidelinks,bookmarks=false]{hyperref}
\newcommand{\ie}{i.e.\ }
\newcommand{\cf}{cf.\ }
\newcommand{\resp}{respectively\ }
\newcommand{\Subsection}{\subsection}
\providecommand{\nobreakdash}{\nobreak\hbox{-}}
\setlength{\emergencystretch}{2em}
\multlinegap0pt
\usepackage{amssymb}
\usepackage{float}
\usepackage{tcolorbox}
\usepackage{enumerate}
\usepackage{xcolor}
\usepackage[normalem]{ulem}
\usepackage{algorithm}\usepackage{algpseudocode}
\newtheorem{lemma}{Lemma}
\title{Extended SQP Methods in Nonsmooth Difference Programming Applied to Problems with Variational Inequality Constraints}
\author{Boris S. Mordukhovich, Yixia Song, Shangzhi Zeng, Jin Zhang}
\date{Source: arXiv:2504.05609v3 (CC BY 4.0)}
\begin{document}
\maketitle

\noindent\emph{Source and licence.} Extended SQP Methods in Nonsmooth Difference Programming Applied to Problems with Variational Inequality Constraints. Version arXiv:2504.05609v3 (CC BY 4.0), distributed by arXiv under the Creative Commons Attribution 4.0 International licence, http://creativecommons.org/licenses/by/4.0/, as recorded in the arXiv licence metadata for this exact version and shown on the abstract page https://arxiv.org/abs/2504.05609v3. Full licence notice: \texttt{licenses/arxiv-2504.05609-CC-BY-4.0.txt}.\par\smallskip

\noindent\textit{Editorial scope.} Counted unit: the article's lemma lemma:descent\_direction (source0.tex lines 551--558). The article's complete original proof of it is delivered with it (source0.tex lines 561--604, one \texttt{proof} environment of 44 lines). No mathematics is written, repaired or re-derived: the statement and the proof are the article's own lines, in the article's own notation, and no retained line is substituted or re-flowed.  Retained uncounted as the source-local context the proof needs: lines 197--199; lines 246--262; lines 290--297; lines 347--351; lines 447--456; lines 449--451; lines 453--455; lines 460--465.  Nothing was omitted from any retained span, so the package carries no omission record.  No retained line contains a citation, so the wrapper carries no bibliography and no bibliography entry.  No retained line contains a figure environment, an image-inclusion command or drawing code, so no image is delivered and the package declares no asset dependency.  Editorial formatting only, in addition: an AMS \texttt{amsart} preamble that restates the article's own declarations and macros used by the retained lines, namely the environments \texttt{lemma} on separate counters, together with the standard packages those lines need.  The retained environments print in this document as Lemma 1 in order of appearance, whereas the article prints the same items with the numbering of its own full document.  The complete native source of the article as distributed in the arXiv e-print for this exact version is bundled unchanged as \texttt{sources/176154/source0.tex}.
\begin{equation}\label{nc-conv}
\mathcal{N}_D(x) := \big\{ v \in \mathbb{R}^n ~\big|~ \langle v, x' - x\rangle \le 0\;\mbox{  for all }\;x' \in D\big\}
\end{equation}

\begin{lemma}\label{max rule}
Let \( f_i: \mathbb{R}^n \to (-\infty,\infty ] \), \( i \in \mathcal{I} \), be a finite collection of proper functions, and let \( \bar{x} \in \bigcap_{i \in \mathcal{I}} \text{\rm dom}(f_i) \). Suppose that \( f_i \) is locally Lipschitzian around this point for all \( i \in \mathcal{I} \). Then the maximum function  
\( f = \max_{i \in \mathcal{I}} f_i \) is also locally Lipschitzian around \( \bar{x} \), and for any \( d \in \mathbb{R}^n \) we have 
\begin{equation}\label{max rule formula}
\begin{aligned}
f'(\bar{x}; d) = \max_{i \in \mathcal{I}_{\max}(\bar{x})} \left\{ (f_i)'(\bar{x}; d) \right\} \leq 
\max_{i \in \mathcal{I}} \left\{ f_i(\bar{x}) + (f_i)'(\bar{x}; d) \right\} - f(\bar{x}),
\end{aligned}
\end{equation}  
where \( \mathcal{I}_{\max}(\bar{x}) := \left\{ i \in \mathcal{I} \mid f_i(\bar{x}) = f(\bar{x}) \right\} \).  
In particular, it holds for \( \mathcal{I} = \{1, 2\} \) and \( f_2 \equiv 0 \) that
\begin{equation*}
\begin{aligned}
\left( \max\{f_1, 0\} \right)'(\bar{x}; d) \leq \max\{ f_1(\bar x) + (f_1)'(\bar{x}; d), 0 \} - \max\{ f_1(\bar x), 0 \}.
\end{aligned}
\end{equation*}
\end{lemma}

\begin{lemma}\label{minus prox regular}
Let $f:\mathbb{R}^n\to(-\infty,\infty]$ be locally Lipschitzian around $\bar x$ and prox-regular at this point. Then $f$ is lower regular at $\bar x$ with the property $\mathrm{co}\,\partial(-f)(\bar x) = -\partial f(\bar x)$ and the representation
\begin{equation*}
\begin{aligned}
(-f)'(\bar x,d) = \inf\left\{\langle -v,d \rangle \mid v \in \partial f(\bar x)\right\}\;\mbox{ for any }\;d\in\mathbb{R}^n.
\end{aligned}
\end{equation*}
\end{lemma}

\begin{equation}\label{direction_sub}
\begin{aligned}
\min_{d \in \mathbb{R}^n} \, \tilde{ \phi}_{k}(x^k +d) + {\frac{1}{2} \langle G_kd,d\rangle} + \delta_X(x^k + d),
\end{aligned}
\end{equation}

\begin{lemma}\label{lem:dk_foc}
Let $d^k$ be the solution to subproblem \eqref{direction_sub}. Then there {exist $\lambda_{k,i} \in [0,1]$ for $i \in \mathcal{I}$} such that
\begin{equation}\label{eq:foc} 
0 \in \frac{1}{p_k}\left(\nabla g_0(x^k)-v_0^k\right) + {\sum_{i\in\mathcal{I}} \lambda_{k,i} \left(\nabla g_i(x^k)-v_i^k\right)} + {G_k} d^k+\mathcal N_X \left(x^k+d^k\right), 
\end{equation}
and, for each $i\in\mathcal{I}$, the following conditions hold
\begin{equation}\label{eq:multiplier}
\lambda_{k,i} \min\left\{ \varphi_i(x^k) + \langle \nabla g_i(x^k) - v_i^k, d^k \rangle, 0 \right\}=(1- \lambda_{k,i})\max\left\{ \varphi_i(x^k) +  \langle \nabla g_i(x^k) - v_i^k, d^k \rangle, 0 \right\} =0.
\end{equation} 
\end{lemma}\vspace*{0.05in}

\begin{equation} 
0 \in \frac{1}{p_k}\left(\nabla g_0(x^k)-v_0^k\right) + {\sum_{i\in\mathcal{I}} \lambda_{k,i} \left(\nabla g_i(x^k)-v_i^k\right)} + {G_k} d^k+\mathcal N_X \left(x^k+d^k\right), 
\end{equation}

\begin{equation}
\lambda_{k,i} \min\left\{ \varphi_i(x^k) + \langle \nabla g_i(x^k) - v_i^k, d^k \rangle, 0 \right\}=(1- \lambda_{k,i})\max\left\{ \varphi_i(x^k) +  \langle \nabla g_i(x^k) - v_i^k, d^k \rangle, 0 \right\} =0.
\end{equation} 

\begin{lemma}\label{lemma:multiplier_inequality}
Let $d^k$ solve subproblem \eqref{direction_sub}. Then for any {$\lambda_{k,i} \in \left[0, 1\right]$ ($i\in\mathcal{I}$)} satisfying \eqref{eq:multiplier}, we have, for each {$i \in \mathcal{I}$}
\begin{equation}\label{eq:multiplier_inequality}
\max\left\{\varphi_i(x^k) + \langle \nabla g_i(x^k) - v_i^k, d^k \rangle, 0\right\} - \max\left\{\varphi_i(x^k), 0\right\} \le  \lambda_{k,i} \langle \nabla g_i(x^k) - v_i^k, d^k \rangle.
\end{equation}
\end{lemma}

\begin{lemma}\label{lemma:descent_direction}
Let the direction $d^k$ solve subproblem \eqref{direction_sub}. Then there exist {$\lambda_{k,i} \in [0,1]$ for $i\in\mathcal{I}$} satisfying  conditions \eqref{eq:foc}, \eqref{eq:multiplier} and such that
\begin{equation}\label{eq:descent_direction}
(\phi_{p_k})'(x^k; d^k) \le 
\frac{1}{p_k} \left\langle \nabla g_0(x^k) - v_0^k, d^k \right\rangle 
+ \sum_{i\in\mathcal I} \lambda_{k,i} \left\langle\nabla g_i(x^k)-v_i^k, d^k\right\rangle \le -\langle G_kd^k,d^k \rangle.
\end{equation}
\end{lemma}

\begin{proof}\,
Since $d^k$ solves the convex subproblem \eqref{direction_sub}, Lemma~\ref{lem:dk_foc} provides the existence of $\lambda_{k,i} \in [0,1]$ ($i \in \mathcal{I}$) satisfying \eqref{eq:foc} and \eqref{eq:multiplier}. 
By definition of the upper directional derivative, we have
\begin{equation}\label{prop1_eq1}
  (\phi_{p_k})'(x^k; d^k) 
  \leq \frac{1}{p_k} \left( \langle \nabla g_0(x^k), d^k \rangle + (-h_0)'(x^k; d^k)\right)
  + \sum_{i\in\mathcal I} \left(\max\{\varphi_i, 0\}\right)'(x^k; d^k).
\end{equation}
Remembering that $h_0$ is locally Lipschitzian around $x^k$ and prox-regular at $x^k$, and that $v_0^k \in \partial h_0(x^k)$, we get by Lemma~\ref{minus prox regular} that
\begin{equation}\label{prop1_eq2}
 (-h_0)'(x^k; d^k) = \inf\left\{ \langle -v, d^k \rangle \mid v \in \partial h_0(x^k) \right\} \leq \langle -v_0^k, d^k \rangle.
\end{equation}
Applying now Lemma~\ref{max rule} to the terms $\left(\max\{\varphi_i, 0\}\right)'(x^k; d^k)$ tells us that for each $i \in \mathcal I$,
\begin{equation}\label{prop1_eq3}
 \left(\max\{{\varphi_i}, 0\}\right)'\left(x ^k; d^k\right) \le \max\left\{ {\varphi_i}(x^k) +{\varphi_i}'(x^k; d^k), 0 \right\} - \max\left\{ {\varphi_i}(x^k), 0 \right\}.
\end{equation}
Since $h_i$ is also locally Lipschitzian around $x^k$ and prox-regular at this point with $v_i^k \in \partial h_i(x^k)$, it follows from Lemma~\ref{minus prox regular} that
\begin{equation*}
 (-h_i)'(x^k; d^k)  \leq \langle -v_i^k, d^k \rangle.
\end{equation*}
Combining the latter with the smoothness of {$g_i$} yields for each $i \in \mathcal I$
 \begin{equation*}
   \varphi_i'(x^k; d^k) = \langle \nabla {g_i}(x^k), d^k \rangle + (-h_i)'(x^k; d^k) \le \langle \nabla g_i(x^k) - v_i^k, d^k \rangle,
 \end{equation*}
and therefore, summing over $i\in\mathcal I$, we deduce from \eqref{prop1_eq3} that
\begin{equation}\label{prop1_eq4}
 \begin{aligned}
   \sum_{i\in\mathcal I} \left(\max\{\varphi_i, 0\}\right)'\left(x^k; d^k\right) \leq\, &\sum_{i\in\mathcal I} \left( \max\left\{ \varphi_i(x^k) + \langle \nabla g_i(x^k) - v_i^k, d^k \rangle, 0 \right\} 
   - \max\left\{ \varphi_i(x^k), 0 \right\} \right) \\
   \le \, & \sum_{i\in\mathcal I} \lambda_{k,i} \left\langle \nabla g_i(x^k) - v_i^k, d^k \right\rangle,
 \end{aligned}
\end{equation}
while the second inequality therein follows from Lemma~\ref{lemma:multiplier_inequality}. Combining \eqref{prop1_eq1}, \eqref{prop1_eq2} and \eqref{prop1_eq4} yields the first estimate in \eqref{eq:descent_direction}. Next we use \eqref{eq:foc}, 
the fact that $x^k, x^k + d^k \in X$, and definition \eqref{nc-conv} of the normal cone $\mathcal{N}_X$ to obtain
\begin{equation*}
\left\langle \frac{1}{p_k}\left(\nabla g_0(x^k)-v_0^k\right) + \sum_{i\in\mathcal I} \lambda_{k,i} \left(\nabla g_i(x^k)-v_i^k\right) + G_kd^k, x^k + d^k - x^k \right \rangle \le 0,
\end{equation*}
which provides in turn the estimate
\begin{equation}\label{prop1_eq5}
\frac{1}{p_k}\left\langle\nabla g_0(x^k)-v_0^k, d^k\right\rangle + \sum_{i\in\mathcal I} \lambda_{k,i} \left\langle\nabla g_i(x^k)-v_i^k, d^k\right\rangle \le -\langle G_kd^k,d^k \rangle.
\end{equation}
%Combining \eqref{prop1_eq1}, \eqref{prop1_eq2}, \eqref{prop1_eq4}, and
Then \eqref{prop1_eq5} verifies the claimed inequalities in \eqref{eq:descent_direction}.
\end{proof}\vspace*{0.05in}

\end{document}
